% ====================================================================

% ====================================================================
\documentclass[article,paper=a4paper]{jlreq}

% jpreportパッケージの読み込み
\usepackage{jpreport}

% ページレイアウトの設定
\usepackage{geometry}
\usepackage{enumitem}
\geometry{
    top=25mm,
    bottom=25mm,
    left=25mm,
    right=25mm
}
% TikZライブラリ
\usetikzlibrary{
    patterns,
    calc,
    arrows,
    decorations,
    angles,
    shapes.geometric,
    fit,
    knots
}
% TikZ用のスタイル
\tikzset{
  every path/.style={
    black,
    line width=1pt
  },
  every node/.style={
    transform shape,
    knot crossing,
    inner sep=1.5pt
  }
}
% ハイパーリンクの設定
\usepackage[unicode]{hyperref}
\usepackage{bookmark}
\usepackage[nameinlink]{cleveref}

\hypersetup{
    colorlinks=true,
    linkcolor=black,
    urlcolor=blue,
    citecolor=blue,
    pdftitle={制御理論入門},
    pdfauthor={akari0koutya}
}

% 数学演算子の定義


% cleveref用の日本語設定
\crefname{thm}{定理}{定理}
\crefname{lem}{補題}{補題}
\crefname{cor}{系}{系}
\crefname{ex}{例}{例}
\crefname{def}{定義}{定義}
\crefname{dfn}{定義}{定義}
\crefname{prop}{命題}{命題}
\crefname{prob}{問題}{問題}
\crefname{rem}{注意}{注意}
\crefname{exe}{演習}{演習}

% 色彩の定義
\definecolor{yukari}{RGB}{196, 163, 191}

% 文書情報
\title{制御概論}
\author{akari0koutya}
\date{\today}

\begin{document}

% --- タイトル ------------------------------------------------------
\maketitle
\tableofcontents
%% 本文の内容をここに記述
% ====================================================================

\clearpage

\section{状態方程式（1）(2026/04/17)}

ここでは，入力をもたない状態方程式
\[
  \dot{x}=Ax
\]
を考える．初期値を $x(0)$ とすると，解は行列の指数関数を用いて
\[
  x(t)=e^{tA}x(0)
\]
と表される．したがって，この節の主な課題は $e^{tA}$ を具体的に計算することである．
これは，スカラーの微分方程式 $\dot{x}=ax$ の解が $x(t)=e^{at}x(0)$ であることの
行列版である．ただし，行列の場合には $A$ の固有値や標準形によって
$e^{tA}$ の見た目が変わるため，固有値の型に分けて計算する．

\subsection{行列指数関数と標準形}

\begin{dfn}
正方行列 $A$ に対して，行列指数関数を
\[
  e^{tA}
    =
    E+tA+\frac{t^2}{2!}A^2+\frac{t^3}{3!}A^3+\cdots
    =
    \sum_{k=0}^{\infty}\frac{t^k}{k!}A^k
\]
で定める．ここで $E$ は単位行列である．
\end{dfn}

\begin{rem}
$e^{tA}$ は数の指数関数と同じ形の級数で定義される．ただし，$A^2,A^3,\ldots$ は行列の積である．
この定義により，$A$ が対角行列に近い形で表せるときには，各成分ごとの指数関数として扱える．
\end{rem}

\begin{prop}
状態方程式
\[
  \dot{x}=Ax
\]
の初期値 $x(0)$ に対する解は
\[
  x(t)=e^{tA}x(0)
\]
である．
\end{prop}

\begin{proof}
行列指数関数を項別に微分する．
\begin{align*}
  \frac{d}{dt}e^{tA}
    &=
    \frac{d}{dt}
    \left(
      E+tA+\frac{t^2}{2!}A^2+\frac{t^3}{3!}A^3+\cdots
    \right)\\
    &=
    A+tA^2+\frac{t^2}{2!}A^3+\cdots\\
    &=
    A\left(E+tA+\frac{t^2}{2!}A^2+\cdots\right)\\
    &=Ae^{tA}.
\end{align*}
したがって，$x(t)=e^{tA}x(0)$ とおけば
\begin{align*}
  \dot{x}(t)
    &=
    \frac{d}{dt}\{e^{tA}x(0)\}\\
    &=
    \left(\frac{d}{dt}e^{tA}\right)x(0)\\
    &=
    Ae^{tA}x(0)
     =
    Ax(t).
\end{align*}
つまり，$x(t)=e^{tA}x(0)$ は $\dot{x}=Ax$ を満たす．また
\[
  x(0)=e^{0A}x(0)=Ex(0)=x(0)
\]
であるから，これは与えられた初期値を満たす解である．
\end{proof}

\begin{prop}
$A$ が正則行列 $P$ により
\[
  P^{-1}AP=J
\]
と変換できるとき，
\[
  e^{tA}=Pe^{tJ}P^{-1}
\]
である．
\end{prop}

\begin{proof}
$P^{-1}AP=J$ より $A=PJP^{-1}$ である．例えば
\[
  A^2
    =
    (PJP^{-1})(PJP^{-1})
    =
    PJ(P^{-1}P)JP^{-1}
    =
    PJ^2P^{-1}
\]
であり，同様に
\[
  A^3
    =
    A^2A
    =
    (PJ^2P^{-1})(PJP^{-1})
    =
    PJ^3P^{-1}
\]
である．一般に，$k=0,1,2,\ldots$ に対して
\[
  A^k=(PJP^{-1})^k=P J^k P^{-1}
\]
が成り立つ．これを行列指数関数の級数に代入すると
\begin{align*}
  e^{tA}
    &=
    \sum_{k=0}^{\infty}\frac{t^k}{k!}A^k\\
    &=
    \sum_{k=0}^{\infty}\frac{t^k}{k!}P J^k P^{-1}\\
    &=
    P\left(\sum_{k=0}^{\infty}\frac{t^k}{k!}J^k\right)P^{-1}
     =
    Pe^{tJ}P^{-1}.
\end{align*}
\end{proof}

\subsection{2次行列の場合}

$A\in M_2(\R)$ の場合，固有値の形によって $e^{tA}$ の公式は次の3種類に分かれる．
実固有値が2つに分かれる場合，実固有値が重なる場合，複素共役の固有値をもつ場合である．
以下の公式を使えば，固有ベクトルを明示的に並べた行列 $P$ を最後まで計算しなくても，
$A$ と固有値だけから $e^{tA}$ を求められる．

\begin{prop}
\textbf{相異なる実固有値の場合．}
$A$ の固有値が $\lambda_1,\lambda_2\in\R$，$\lambda_1\ne\lambda_2$ であるとする．このとき
\[
  e^{tA}
    =
    e^{\lambda_1t}P_1+e^{\lambda_2t}P_2
\]
である．ただし
\[
  P_1=\frac{1}{\lambda_1-\lambda_2}(A-\lambda_2E),
  \qquad
  P_2=\frac{1}{\lambda_2-\lambda_1}(A-\lambda_1E)
\]
である．
\end{prop}

\begin{proof}
固有ベクトルを $p_1,p_2$ とし，
\[
  P=(p_1\ p_2)
\]
とおくと
\[
  P^{-1}AP=
  \begin{pmatrix}
    \lambda_1 & 0\\
    0 & \lambda_2
  \end{pmatrix}.
\]
対角行列の指数関数は対角成分ごとに指数関数をとればよいので
\[
  e^{t(P^{-1}AP)}
  =
  \begin{pmatrix}
    e^{\lambda_1t} & 0\\
    0 & e^{\lambda_2t}
  \end{pmatrix}
\]
である．したがって，相似変換の公式より
\[
  e^{tA}
    =
    P
    \begin{pmatrix}
      e^{\lambda_1t} & 0\\
      0 & e^{\lambda_2t}
    \end{pmatrix}
    P^{-1}.
\]
ここで
\[
  E_1=
  \begin{pmatrix}
    1 & 0\\
    0 & 0
  \end{pmatrix},
  \qquad
  E_2=
  \begin{pmatrix}
    0 & 0\\
    0 & 1
  \end{pmatrix}
\]
とおくと
\[
  P^{-1}AP=\lambda_1E_1+\lambda_2E_2,
  \qquad
  E_1+E_2=E.
\]
両辺を左から $P$，右から $P^{-1}$ で挟むと
\[
  A=\lambda_1PE_1P^{-1}+\lambda_2PE_2P^{-1}
\]
である．$E_2=E-E_1$ を用いると
\begin{align*}
  A
    &=
    \lambda_1PE_1P^{-1}
    +\lambda_2P(E-E_1)P^{-1}\\
    &=
    (\lambda_1-\lambda_2)PE_1P^{-1}+\lambda_2E.
\end{align*}
したがって
\[
  PE_1P^{-1}
    =
    \frac{1}{\lambda_1-\lambda_2}(A-\lambda_2E).
\]
同様に
\[
  E_1=E-E_2
\]
を用いれば
\begin{align*}
  A
    &=
    \lambda_1P(E-E_2)P^{-1}
    +\lambda_2PE_2P^{-1}\\
    &=
    \lambda_1E+(\lambda_2-\lambda_1)PE_2P^{-1}
\end{align*}
である．よって
\[
  PE_2P^{-1}
    =
    \frac{1}{\lambda_2-\lambda_1}(A-\lambda_1E).
\]
これらを $P_1,P_2$ と書けば，
\[
  e^{tA}
    =
    e^{\lambda_1t}P_1+e^{\lambda_2t}P_2
\]
を得る．
\end{proof}

\begin{prop}
\textbf{重解の場合．}
$A$ の固有値が $\lambda$ の重解で，ジョルダン標準形が
\[
  P^{-1}AP=
  \begin{pmatrix}
    \lambda & 1\\
    0 & \lambda
  \end{pmatrix}
\]
であるとする．このとき
\[
  e^{tA}
    =
    e^{\lambda t}\{E+t(A-\lambda E)\}
\]
である．
\end{prop}

\begin{proof}
\[
  D=
  \begin{pmatrix}
    0 & 1\\
    0 & 0
  \end{pmatrix}
\]
とおくと，$D^2=O$ であり，
\[
  P^{-1}AP=\lambda E+D
\]
と書ける．よって
\[
  A=\lambda E+PDP^{-1}
\]
である．また
\[
  e^{tD}
    =
    E+tD+\frac{t^2}{2!}D^2+\cdots
    =
    E+tD
\]
である．第2項以降で $D^2$ を含むものはすべて零行列になるからである．したがって
\begin{align*}
  e^{tA}
    &=
    P e^{t(\lambda E+D)} P^{-1}\\
    &=
    e^{\lambda t}P(E+tD)P^{-1}\\
    &=
    e^{\lambda t}\{E+tPDP^{-1}\}.
\end{align*}
ここで $PDP^{-1}=A-\lambda E$ であるから，
\[
  e^{tA}
    =
    e^{\lambda t}\{E+t(A-\lambda E)\}
\]
を得る．
\end{proof}

\begin{prop}
\textbf{複素共役固有値の場合．}
$A$ の固有値が $\alpha\pm i\beta$，$\beta\ne0$ であるとする．このとき
\[
  e^{tA}
    =
    e^{\alpha t}
    \left\{
      \cos(\beta t)E
      +\frac{\sin(\beta t)}{\beta}(A-\alpha E)
    \right\}
\]
である．
\end{prop}

\begin{proof}
実標準形を
\[
  P^{-1}AP=
  \begin{pmatrix}
    \alpha & \beta\\
    -\beta & \alpha
  \end{pmatrix}
  =
  \alpha E+\beta J,
  \qquad
  J=
  \begin{pmatrix}
    0 & 1\\
    -1 & 0
  \end{pmatrix}
\]
と書く．このとき $J^2=-E$ であるから
\begin{align*}
  e^{\beta tJ}
    &=
    E+\beta tJ+\frac{(\beta t)^2}{2!}J^2
      +\frac{(\beta t)^3}{3!}J^3+\cdots\\
    &=
    \left(1-\frac{(\beta t)^2}{2!}
      +\frac{(\beta t)^4}{4!}-\cdots\right)E\\
    &\quad+
    \left(\beta t-\frac{(\beta t)^3}{3!}
      +\frac{(\beta t)^5}{5!}-\cdots\right)J\\
    &=
    \cos(\beta t)E+\sin(\beta t)J.
\end{align*}
したがって
\begin{align*}
  e^{tA}
    &=
    P e^{t(\alpha E+\beta J)}P^{-1}\\
    &=
    e^{\alpha t}
    P\{\cos(\beta t)E+\sin(\beta t)J\}P^{-1}\\
    &=
    e^{\alpha t}
    \{\cos(\beta t)E+\sin(\beta t)PJP^{-1}\}.
\end{align*}
一方，
\[
  A=P(\alpha E+\beta J)P^{-1}
    =
    \alpha E+\beta PJP^{-1}
\]
なので
\[
  PJP^{-1}=\frac{1}{\beta}(A-\alpha E).
\]
これを代入して公式を得る．
\end{proof}

\begin{rem}
$2$ 次行列
\[
  A=
  \begin{pmatrix}
    a & b\\
    c & d
  \end{pmatrix}
\]
の固有値は
\[
  \det(\lambda E-A)
    =
    \begin{vmatrix}
      \lambda-a & -b\\
      -c & \lambda-d
    \end{vmatrix}
    =
    \lambda^2-(\operatorname{tr}A)\lambda+\det A
\]
から求められる．すなわち
\[
  \lambda^2-(\operatorname{tr}A)\lambda+\det A=0
\]
を解けばよい．
\end{rem}

\subsection{例}

\begin{ex}
次の微分方程式の解を求める．ただし，初期値を $x(0)$ とする．また，
$t\to\infty$ での解の挙動を調べる．
\[
  \begin{array}{lll}
    \text{(1)}&
    \displaystyle
    \dot{x}=
    \begin{pmatrix}
      7 & -6\\
      3 & -2
    \end{pmatrix}x,
    &
    \text{(2)}\quad 
    \displaystyle
    \dot{x}=
    \begin{pmatrix}
      3 & -1\\
      1 & 1
    \end{pmatrix}x,\\[5mm]
    \text{(3)}&
    \displaystyle
    \dot{x}=
    \begin{pmatrix}
      -1 & 2\\
      -1 & -3
    \end{pmatrix}x.
  \end{array}
\]
\end{ex}

\begin{proof}[解]
\textbf{(1)}
\[
  A=
  \begin{pmatrix}
    7 & -6\\
    3 & -2
  \end{pmatrix}
\]
とおく．
まず固有値を調べ，どの公式を使うかを決める．
\[
  \operatorname{tr}A=7+(-2)=5,
  \qquad
  \det A=7\cdot(-2)-(-6)\cdot3=4
\]
であるから，特性方程式は
\[
  \lambda^2-5\lambda+4=0
  \qquad\Longleftrightarrow\qquad
  (\lambda-1)(\lambda-4)=0.
\]
よって固有値は $\lambda_1=1,\lambda_2=4$ である．相異なる実固有値の場合の公式より
\[
  e^{tA}=e^{t}P_1+e^{4t}P_2
\]
となる．ここで
\begin{align*}
  P_1
    &=
    \frac{1}{1-4}(A-4E)
     =
    -\frac{1}{3}
    \left\{
    \begin{pmatrix}
      7 & -6\\
      3 & -2
    \end{pmatrix}
    -
    \begin{pmatrix}
      4 & 0\\
      0 & 4
    \end{pmatrix}
    \right\}\\
    &=
    -\frac{1}{3}
    \begin{pmatrix}
      3 & -6\\
      3 & -6
    \end{pmatrix}
     =
    \begin{pmatrix}
      -1 & 2\\
      -1 & 2
    \end{pmatrix},\\
  P_2
    &=
    \frac{1}{4-1}(A-E)
     =
    \frac{1}{3}
    \left\{
    \begin{pmatrix}
      7 & -6\\
      3 & -2
    \end{pmatrix}
    -
    \begin{pmatrix}
      1 & 0\\
      0 & 1
    \end{pmatrix}
    \right\}\\
    &=
    \frac{1}{3}
    \begin{pmatrix}
      6 & -6\\
      3 & -3
    \end{pmatrix}
     =
    \begin{pmatrix}
      2 & -2\\
      1 & -1
    \end{pmatrix}.
\end{align*}
したがって
\[
  e^{tA}
    =
    e^t
    \begin{pmatrix}
      -1 & 2\\
      -1 & 2
    \end{pmatrix}
    +
    e^{4t}
    \begin{pmatrix}
      2 & -2\\
      1 & -1
    \end{pmatrix}.
\]
よって
\[
  x(t)
    =
    \left\{
      e^t
      \begin{pmatrix}
        -1 & 2\\
        -1 & 2
      \end{pmatrix}
      +
      e^{4t}
      \begin{pmatrix}
        2 & -2\\
        1 & -1
      \end{pmatrix}
    \right\}x(0).
\]
この解は，$e^t$ を含む成分と $e^{4t}$ を含む成分の和である．
どちらも $t\to\infty$ で増大するため，零解以外では原点から離れていく．
したがって $x(0)\ne0$ なら解は発散し，
\[
  \|x(t)\|\to\infty
  \qquad (t\to\infty)
\]
となる．

\medskip
\textbf{(2)}
\[
  A=
  \begin{pmatrix}
    3 & -1\\
    1 & 1
  \end{pmatrix}
\]
とおく．
ここでもまず特性方程式を作る．
\[
  \operatorname{tr}A=3+1=4,
  \qquad
  \det A=3\cdot1-(-1)\cdot1=4
\]
であるから，
\[
  \lambda^2-4\lambda+4=0
  \qquad\Longleftrightarrow\qquad
  (\lambda-2)^2=0.
\]
したがって固有値は $\lambda=2$ の重解である．
この場合には重解の公式
\[
  e^{tA}=e^{2t}\{E+t(A-2E)\}
\]
を用いる．実際，
\[
  A-2E=
  \begin{pmatrix}
    3 & -1\\
    1 & 1
  \end{pmatrix}
  -
  \begin{pmatrix}
    2 & 0\\
    0 & 2
  \end{pmatrix}
  =
  \begin{pmatrix}
    1 & -1\\
    1 & -1
  \end{pmatrix}
\]
なので
\[
  e^{tA}
    =
    e^{2t}\{E+t(A-2E)\}
    =
    e^{2t}
    \left\{
      \begin{pmatrix}
        1 & 0\\
        0 & 1
      \end{pmatrix}
      +
      t
      \begin{pmatrix}
        1 & -1\\
        1 & -1
      \end{pmatrix}
    \right\}.
\]
よって
\[
  x(t)
    =
    e^{2t}
    \left\{
      \begin{pmatrix}
        1 & 0\\
        0 & 1
      \end{pmatrix}
      +
      t
      \begin{pmatrix}
        1 & -1\\
        1 & -1
      \end{pmatrix}
    \right\}x(0).
\]
括弧の中には $t$ に比例する項があるが，支配的なのは前に掛かっている指数因子 $e^{2t}$ である．
$e^{2t}$ が増大するため，$x(0)\ne0$ なら
\[
  \|x(t)\|\to\infty
  \qquad (t\to\infty)
\]
である．

\medskip
\textbf{(3)}
\[
  A=
  \begin{pmatrix}
    -1 & 2\\
    -1 & -3
  \end{pmatrix}
\]
とおく．
固有値を求めると，
\[
  \operatorname{tr}A=-1+(-3)=-4,
  \qquad
  \det A=(-1)(-3)-2(-1)=5
\]
であるから，
\[
  \lambda^2+4\lambda+5=0
  \qquad\Longleftrightarrow\qquad
  (\lambda+2)^2+1=0.
\]
したがって固有値は
\[
  \lambda=-2\pm i
\]
である．ここでは $\alpha=-2,\beta=1$ であり，
\[
  A-\alpha E=A+2E=
  \begin{pmatrix}
    -1 & 2\\
    -1 & -3
  \end{pmatrix}
  +
  \begin{pmatrix}
    2 & 0\\
    0 & 2
  \end{pmatrix}
  =
  \begin{pmatrix}
    1 & 2\\
    -1 & -1
  \end{pmatrix}.
\]
複素共役固有値の場合の公式より
\[
  e^{tA}
    =
    e^{-2t}
    \left\{
      \cos(t)E+\sin(t)(A+2E)
    \right\}.
\]
すなわち
\[
  x(t)
    =
    e^{-2t}
    \left\{
      \cos(t)
      \begin{pmatrix}
        1 & 0\\
        0 & 1
      \end{pmatrix}
      +
      \sin(t)
      \begin{pmatrix}
        1 & 2\\
        -1 & -1
      \end{pmatrix}
    \right\}x(0).
\]
この解は，回転を表す三角関数の部分に，減衰因子 $e^{-2t}$ が掛かった形である．
三角関数の部分は有界であり，$e^{-2t}\to0$ であるから
\[
  \|x(t)\|\to0
  \qquad (t\to\infty)
\]
である．
\end{proof}

\begin{exe}
次の同次状態方程式についても，同様に $e^{tA}$ を求めて解の挙動を調べよ．
\[
  \begin{array}{lll}
    \text{(1)}&
    \displaystyle
    \dot{x}=
    \begin{pmatrix}
      -5 & 4\\
      -2 & 1
    \end{pmatrix}x,
    &
    \text{(2)}\quad 
    \displaystyle
    \dot{x}=
    \begin{pmatrix}
      3 & -4\\
      1 & -1
    \end{pmatrix}x,\\[5mm]
    \text{(3)}&
    \displaystyle
    \dot{x}=
    \begin{pmatrix}
      1 & -4\\
      2 & -3
    \end{pmatrix}x.
  \end{array}
\]
\end{exe}

\begin{proof}[解答]
\textbf{(1)}
\[
  A=
  \begin{pmatrix}
    -5 & 4\\
    -2 & 1
  \end{pmatrix}
\]
とおく．まず固有値を求める．
\[
  \operatorname{tr}A=-5+1=-4,
  \qquad
  \det A=(-5)\cdot1-4\cdot(-2)=3
\]
であるから，特性方程式は
\[
  \lambda^2+4\lambda+3=0
  \qquad\Longleftrightarrow\qquad
  (\lambda+1)(\lambda+3)=0.
\]
したがって固有値は $\lambda_1=-1,\lambda_2=-3$ である．
相異なる実固有値の場合の公式を用いると
\[
  e^{tA}=e^{-t}P_1+e^{-3t}P_2
\]
である．ここで
\begin{align*}
  P_1
    &=
    \frac{1}{-1-(-3)}\{A-(-3)E\}
     =
    \frac{1}{2}
    \left\{
    \begin{pmatrix}
      -5 & 4\\
      -2 & 1
    \end{pmatrix}
    +
    \begin{pmatrix}
      3 & 0\\
      0 & 3
    \end{pmatrix}
    \right\} \\
    &=
    \frac{1}{2}
    \begin{pmatrix}
      -2 & 4\\
      -2 & 4
    \end{pmatrix}
     =
    \begin{pmatrix}
      -1 & 2\\
      -1 & 2
    \end{pmatrix},\\
  P_2
    &=
    \frac{1}{-3-(-1)}\{A-(-1)E\}
     =
    -\frac{1}{2}
    \left\{
    \begin{pmatrix}
      -5 & 4\\
      -2 & 1
    \end{pmatrix}
    +
    \begin{pmatrix}
      1 & 0\\
      0 & 1
    \end{pmatrix}
    \right\}\\
    &=
    -\frac{1}{2}
    \begin{pmatrix}
      -4 & 4\\
      -2 & 2
    \end{pmatrix}
     =
    \begin{pmatrix}
      2 & -2\\
      1 & -1
    \end{pmatrix}.
\end{align*}
よって
\[
  x(t)
    =
    \left\{
      e^{-t}
      \begin{pmatrix}
        -1 & 2\\
        -1 & 2
      \end{pmatrix}
      +
      e^{-3t}
      \begin{pmatrix}
        2 & -2\\
        1 & -1
      \end{pmatrix}
    \right\}x(0).
\]
指数因子 $e^{-t}$，$e^{-3t}$ はどちらも $0$ に収束するので
\[
  \|x(t)\|\to0
  \qquad (t\to\infty)
\]
である．

\medskip
\textbf{(2)}
\[
  A=
  \begin{pmatrix}
    3 & -4\\
    1 & -1
  \end{pmatrix}
\]
とおく．このとき
\[
  \operatorname{tr}A=3+(-1)=2,
  \qquad
  \det A=3\cdot(-1)-(-4)\cdot1=1
\]
であるから，
\[
  \lambda^2-2\lambda+1=0
  \qquad\Longleftrightarrow\qquad
  (\lambda-1)^2=0.
\]
したがって固有値は $\lambda=1$ の重解である．重解の場合の公式より
\[
  e^{tA}=e^t\{E+t(A-E)\}
\]
である．実際，
\[
  A-E=
  \begin{pmatrix}
    3 & -4\\
    1 & -1
  \end{pmatrix}
  -
  \begin{pmatrix}
    1 & 0\\
    0 & 1
  \end{pmatrix}
  =
  \begin{pmatrix}
    2 & -4\\
    1 & -2
  \end{pmatrix}
\]
なので
\[
  x(t)
    =
    e^t
    \left\{
      \begin{pmatrix}
        1 & 0\\
        0 & 1
      \end{pmatrix}
      +
      t
      \begin{pmatrix}
        2 & -4\\
        1 & -2
      \end{pmatrix}
    \right\}x(0).
\]
ここで $(A-E)^2=O$ であるから，括弧の中は $x(0)$ に対する高々一次の項である．
もし $(A-E)x(0)=0$ なら $x(t)=e^t x(0)$ となり，$(A-E)x(0)\ne0$ なら
$t e^t(A-E)x(0)$ が支配的になる．したがって，$x(0)\ne0$ なら
\[
  \|x(t)\|\to\infty
  \qquad (t\to\infty)
\]
である．

\medskip
\textbf{(3)}
\[
  A=
  \begin{pmatrix}
    1 & -4\\
    2 & -3
  \end{pmatrix}
\]
とおく．固有値を求めると
\[
  \operatorname{tr}A=1+(-3)=-2,
  \qquad
  \det A=1\cdot(-3)-(-4)\cdot2=5
\]
であるから，
\[
  \lambda^2+2\lambda+5=0
  \qquad\Longleftrightarrow\qquad
  (\lambda+1)^2+4=0.
\]
したがって固有値は
\[
  \lambda=-1\pm2i
\]
である．ここでは $\alpha=-1,\beta=2$ であり，
\[
  A-\alpha E=A+E=
  \begin{pmatrix}
    1 & -4\\
    2 & -3
  \end{pmatrix}
  +
  \begin{pmatrix}
    1 & 0\\
    0 & 1
  \end{pmatrix}
  =
  \begin{pmatrix}
    2 & -4\\
    2 & -2
  \end{pmatrix}.
\]
複素共役固有値の場合の公式より
\[
  e^{tA}
    =
    e^{-t}
    \left\{
      \cos(2t)E+\frac{\sin(2t)}{2}(A+E)
    \right\}.
\]
すなわち
\[
  x(t)
    =
    e^{-t}
    \left\{
      \cos(2t)
      \begin{pmatrix}
        1 & 0\\
        0 & 1
      \end{pmatrix}
      +
      \sin(2t)
      \begin{pmatrix}
        1 & -2\\
        1 & -1
      \end{pmatrix}
    \right\}x(0).
\]
三角関数の部分は有界であり，前に掛かっている $e^{-t}$ が $0$ に収束する．
したがって
\[
  \|x(t)\|\to0
  \qquad (t\to\infty)
\]
である．
\end{proof}

\clearpage

\section{状態方程式（2）(2026/04/24)}

今回は，状態方程式の解を行列の指数関数を用いて表す方法を扱う．
同次方程式の解は $e^{tA}$ で記述され，入力 $u(t)$ がある場合には，
その影響を積分として足し合わせる．後半では，同じ解がラプラス変換からも得られることを確認する．

\subsection{行列の指数関数}

\begin{prop}
$A\in M_2(\R)$ とし，$E$ を単位行列とする．$A$ の固有値は
\[
  \lambda^2-(\operatorname{tr} A)\lambda+\det A=0
\]
の解として求められる．このとき，$e^{tA}$ は固有値の形に応じて次のように表せる．

\begin{enumerate}[label=(\arabic*)]
  \item $\lambda_1\ne\lambda_2$ のとき
  \begin{align*}
    e^{tA}
      &= e^{\lambda_1 t}P_1+e^{\lambda_2 t}P_2,\\
    P_1
      &=\frac{A-\lambda_2 E}{\lambda_1-\lambda_2},
      &
    P_2
      &=\frac{A-\lambda_1 E}{\lambda_2-\lambda_1}.
  \end{align*}

  \item $\lambda=\lambda_0$ が重解のとき
  \[
    e^{tA}
      = e^{\lambda_0 t}\{E+t(A-\lambda_0 E)\}.
  \]

  \item $\lambda=\alpha\pm i\beta$，
  $\alpha\in\R,\ \beta\in\R^\times$ のとき
  \[
    e^{tA}
      = e^{\alpha t}
        \left\{
          \cos(\beta t)E
          +\frac{\sin(\beta t)}{\beta}(A-\alpha E)
        \right\}.
  \]
\end{enumerate}
\end{prop}

\begin{proof}
ここでは公式の意味を確認しておく．$2$ 次正方行列では，ケイリー・ハミルトンの定理により，
$A$ の高い次数のべきは $E$ と $A$ の一次結合に帰着できる．したがって $e^{tA}$ も
$E$ と $A$ の一次結合として表せる．

固有値が異なる場合，$P_1,P_2$ は
\[
  P_1+P_2=E,
  \qquad
  AP_1=\lambda_1P_1,
  \qquad
  AP_2=\lambda_2P_2
\]
を満たす射影行列である．よって
\[
  A=\lambda_1P_1+\lambda_2P_2,
  \qquad
  A^k=\lambda_1^k P_1+\lambda_2^k P_2
  \quad (k=0,1,2,\ldots)
\]
となるから，
\begin{align*}
  e^{tA}
    &=\sum_{k=0}^{\infty}\frac{t^k A^k}{k!} \\
    &=\sum_{k=0}^{\infty}\frac{(\lambda_1t)^k}{k!}P_1
      +\sum_{k=0}^{\infty}\frac{(\lambda_2t)^k}{k!}P_2 \\
    &=e^{\lambda_1t}P_1+e^{\lambda_2t}P_2.
\end{align*}

重解の場合は，$N=A-\lambda_0E$ とおくと $N^2=O$ である．したがって
\[
  e^{tA}
    =e^{t(\lambda_0E+N)}
    =e^{\lambda_0t}e^{tN}
    =e^{\lambda_0t}(E+tN)
    =e^{\lambda_0t}\{E+t(A-\lambda_0E)\}.
\]
複素共役の固有値をもつ場合も，同じく $A-\alpha E$ を用いて $E$ と $A-\alpha E$ に分けると，
三角関数の形の公式が得られる．
\end{proof}

\subsection{解の公式}

\begin{prop}
状態方程式
\begin{equation}
  \dot{x}=Ax+bu(t)
  \label{eq:state-equation}
\end{equation}
の解は
\begin{equation}
  x(t)
    =e^{tA}x(0)
      +\int_0^t e^{(t-\tau)A}bu(\tau)\,d\tau
  \label{eq:solution-formula}
\end{equation}
で与えられる．
\end{prop}

\begin{proof}
まず $u(t)=0$ の場合を考える．このとき
\[
  \dot{x}=Ax
\]
であり，行列指数関数の性質
\[
  \frac{d}{dt}e^{tA}=Ae^{tA}
\]
より，一般解は
\[
  x=e^{tA}c
  \qquad (c:\text{定数})
\]
である．

入力 $u(t)$ がある場合には，定数ベクトル $c$ を時刻に依存する未知関数 $c(t)$ に置き換え，
\[
  x(t)=e^{tA}c(t)
\]
とおく．これを式~\eqref{eq:state-equation} に代入する．まず左辺は積の微分により
\[
  \dot{x}(t)
    =Ae^{tA}c(t)+e^{tA}\dot{c}(t)
\]
である．一方，右辺は
\[
  Ax(t)+bu(t)=Ae^{tA}c(t)+bu(t)
\]
であるから，
\begin{align*}
  Ae^{tA}c(t)+e^{tA}\dot{c}(t)
    &=Ae^{tA}c(t)+bu(t),\\
  e^{tA}\dot{c}(t)
    &=bu(t),\\
  \dot{c}(t)
    &=e^{-tA}bu(t).
\end{align*}
最後の式では，左から $e^{-tA}$ を掛けている．これを $0$ から $t$ まで積分すると
\[
  c(t)-c(0)
    =\int_0^t e^{-\tau A}bu(\tau)\,d\tau
\]
である．ここで $c_0=c(0)$ と書けば
\[
  c(t)=c_0+\int_0^t e^{-\tau A}bu(\tau)\,d\tau.
\]
よって
\begin{align*}
  x(t)
    &=e^{tA}c(t)\\
    &=e^{tA}c_0
      +e^{tA}\int_0^t e^{-\tau A}bu(\tau)\,d\tau\\
    &=e^{tA}c_0
      +\int_0^t e^{tA}e^{-\tau A}bu(\tau)\,d\tau\\
    &=e^{tA}c_0
      +\int_0^t e^{(t-\tau)A}bu(\tau)\,d\tau.
\end{align*}
また，$x(0)=e^{0A}c_0=c_0$ であるから式~\eqref{eq:solution-formula} を得る．
\end{proof}

\begin{rem}
同じ式は，積分因子の考え方からも得られる．
\[
  \dot{x}-Ax=bu(t)
\]
に左から $e^{-tA}$ を掛けると
\[
  e^{-tA}\dot{x}-e^{-tA}Ax=e^{-tA}bu(t)
\]
である．ここで
\[
  \frac{d}{dt}\left(e^{-tA}x(t)\right)
    =
    -Ae^{-tA}x(t)+e^{-tA}\dot{x}(t)
    =
    e^{-tA}\{\dot{x}(t)-Ax(t)\}
\]
だから，
\[
  \frac{d}{dt}\left(e^{-tA}x(t)\right)=e^{-tA}bu(t).
\]
両辺を $0$ から $t$ まで積分すると
\[
  e^{-tA}x(t)-x(0)
    =
    \int_0^t e^{-\tau A}bu(\tau)\,d\tau.
\]
両辺に左から $e^{tA}$ を掛ければ式~\eqref{eq:solution-formula} が従う．
\end{rem}

\subsection{例}

\begin{ex}
次の微分方程式の解を求め，また $t\to\infty$ のときの解の挙動を調べる：
\[
  \begin{cases}
    \displaystyle
    \dot{x}
      =
      \begin{pmatrix}
        1 & -2\\
        3 & -4
      \end{pmatrix}x
      +
      \binom{0}{2}u,\\[3mm]
    u(t)=1.
  \end{cases}
\]
\end{ex}

\begin{proof}[解]
$A=\begin{pmatrix}1&-2\\3&-4\end{pmatrix}$ とおく．
まず，固有値を求める．
\[
  \operatorname{tr}A=1+(-4)=-3,
  \qquad
  \det A=1\cdot(-4)-(-2)\cdot3=2.
\]
したがって特性方程式は
\[
  \lambda^2+3\lambda+2=0
  \qquad\Longleftrightarrow\qquad
  (\lambda+1)(\lambda+2)=0.
\]
よって
\[
  \lambda=-2,\,-1
\]
である．

$\lambda_1=-1,\ \lambda_2=-2$ とすると，
\begin{align*}
  P_1
    &=
    \frac{1}{-1-(-2)}\{A-(-2)E\}
    =
    \begin{pmatrix}
      3 & -2\\
      3 & -2
    \end{pmatrix},\\
  P_2
    &=
    \frac{1}{-2-(-1)}\{A-(-1)E\}
    =
    \begin{pmatrix}
      -2 & 2\\
      -3 & 3
    \end{pmatrix}.
\end{align*}
したがって
\[
  e^{tA}
    =
    e^{-t}
    \begin{pmatrix}
      3 & -2\\
      3 & -2
    \end{pmatrix}
    +
    e^{-2t}
    \begin{pmatrix}
      -2 & 2\\
      -3 & 3
    \end{pmatrix}.
\]

ここで $b=\binom{0}{2}$，$u(t)=1$ より，
\[
  bu(t)=\binom{0}{2}.
\]
入力による項は，$r=t-\tau$ とおくと次のように計算できる．まず
\begin{align*}
  e^{rA}b
    &=
    e^{-r}
    \begin{pmatrix}
      3 & -2\\
      3 & -2
    \end{pmatrix}
    \binom{0}{2}
    +
    e^{-2r}
    \begin{pmatrix}
      -2 & 2\\
      -3 & 3
    \end{pmatrix}
    \binom{0}{2}\\
    &=
    e^{-r}\binom{-4}{-4}
    +
    e^{-2r}\binom{4}{6}.
\end{align*}
したがって
\begin{align*}
  \int_0^t e^{(t-\tau)A}bu(\tau)\,d\tau
    &=
    \int_0^t
      \left\{
        e^{-(t-\tau)}\binom{-4}{-4}
        +
        e^{-2(t-\tau)}\binom{4}{6}
      \right\}d\tau\\
    &=
    \int_0^t
      \left\{
        e^{-r}\binom{-4}{-4}
        +
        e^{-2r}\binom{4}{6}
      \right\}dr\\
    &=
    \binom{-4}{-4}(1-e^{-t})
      +
    \binom{4}{6}\frac{1-e^{-2t}}{2}\\
    &=
    \binom{-2}{-1}
      +
    e^{-t}\binom{4}{4}
      +
    e^{-2t}\binom{-2}{-3}.
\end{align*}
すなわち
\[
  \int_0^t e^{(t-\tau)A}bu(\tau)\,d\tau
    =
    e^{-t}\binom{4}{4}
    +
    e^{-2t}\binom{-2}{-3}
    +
    \binom{-2}{-1}.
\]
よって，解の公式~\eqref{eq:solution-formula} から
\[
  x(t)
    =
    \left\{
      e^{-t}
      \begin{pmatrix}
        3 & -2\\
        3 & -2
      \end{pmatrix}
      +
      e^{-2t}
      \begin{pmatrix}
        -2 & 2\\
        -3 & 3
      \end{pmatrix}
    \right\}x(0)
    +
    e^{-t}\binom{4}{4}
    +
    e^{-2t}\binom{-2}{-3}
    +
    \binom{-2}{-1}.
\]
右辺の $e^{-t}$，$e^{-2t}$ を含む項は，$t\to\infty$ で消えていく過渡項である．
したがって，初期値 $x(0)$ に依存しない定常値として
\[
  \lim_{t\to\infty}x(t)=\binom{-2}{-1}.
\]
\end{proof}

\subsection{ラプラス変換}

\begin{prop}
$x(t)$，$u(t)$ のラプラス変換をそれぞれ $X(s)$，$U(s)$ とする．状態方程式
\[
  \dot{x}=Ax+bu
\]
をラプラス変換すると，
\begin{equation}
  X(s)
    =(sE-A)^{-1}\{bU(s)+x(0)\}
  \label{eq:laplace-solution}
\end{equation}
である．
\end{prop}

\begin{proof}
ラプラス変換では，時間微分が
\[
  \mathcal{L}[\dot{x}](s)=sX(s)-x(0)
\]
に移る．したがって，状態方程式をラプラス変換すると
\begin{align*}
  sX(s)-x(0)
    &=AX(s)+bU(s),\\
  sX(s)-AX(s)
    &=bU(s)+x(0),\\
  (sE-A)X(s)
    &=bU(s)+x(0).
\end{align*}
したがって式~\eqref{eq:laplace-solution} が従う．逆変換により
\[
  x(t)
    =
    \mathcal{L}^{-1}\!\left[(sE-A)^{-1}\right]x(0)
    +
    \mathcal{L}^{-1}\!\left[(sE-A)^{-1}bU(s)\right].
\]
ここで
\[
  \mathcal{L}^{-1}\!\left[(sE-A)^{-1}\right]
    =e^{tA},
\]
であり，積の逆変換には畳み込み定理を用いるので
\[
  \mathcal{L}^{-1}\!\left[(sE-A)^{-1}bU(s)\right]
    =
    \int_0^t e^{(t-\tau)A}bu(\tau)\,d\tau
\]
が成立する．
\end{proof}

\subsection{例のラプラス変換による解法}

\begin{proof}[別解]
先ほどの例をラプラス変換で確認する．
まず
\[
  sE-A
    =
    \begin{pmatrix}
      s-1 & 2\\
      -3 & s+4
    \end{pmatrix}
\]
であり，
\[
  \det(sE-A)
    =(s-1)(s+4)+6
    =s^2+3s+2
    =(s+1)(s+2).
\]
したがって
\[
  (sE-A)^{-1}
    =
    \frac{1}{(s+1)(s+2)}
    \begin{pmatrix}
      s+4 & -2\\
      3 & s-1
    \end{pmatrix}.
\]
これを
\[
  \frac{A_1}{s+1}+\frac{A_2}{s+2}
\]
と部分分数分解する．両辺に $(s+1)(s+2)$ を掛けると
\[
  \begin{pmatrix}
    s+4 & -2\\
    3 & s-1
  \end{pmatrix}
    =
    (s+2)A_1+(s+1)A_2
\]
となる．したがって，$s=-1,-2$ を代入すれば係数が求まる：
\begin{align*}
  A_1
    &=
    \left[
      \frac{1}{s+2}
      \begin{pmatrix}
        s+4 & -2\\
        3 & s-1
      \end{pmatrix}
    \right]_{s=-1}
    =
    \begin{pmatrix}
      3 & -2\\
      3 & -2
    \end{pmatrix},\\
  A_2
    &=
    \left[
      \frac{1}{s+1}
      \begin{pmatrix}
        s+4 & -2\\
        3 & s-1
      \end{pmatrix}
    \right]_{s=-2}
    =
    \begin{pmatrix}
      -2 & 2\\
      -3 & 3
    \end{pmatrix}.
\end{align*}
したがって
\[
  \mathcal{L}^{-1}\!\left[(sE-A)^{-1}\right]
    =
    e^{-t}
    \begin{pmatrix}
      3 & -2\\
      3 & -2
    \end{pmatrix}
    +
    e^{-2t}
    \begin{pmatrix}
      -2 & 2\\
      -3 & 3
    \end{pmatrix}.
\]

また，$b=\binom{0}{2}$，$u(t)=1$ なので $U(s)=1/s$ である．
入力による項は
\[
  (sE-A)^{-1}bU(s)
    =
    \frac{1}{(s+1)(s+2)}
    \begin{pmatrix}
      s+4 & -2\\
      3 & s-1
    \end{pmatrix}
    \binom{0}{2}\frac{1}{s}
    =
    \frac{1}{s(s+1)(s+2)}
    \binom{-4}{2s-2}.
\]
これを
\[
  \frac{B_1}{s}+\frac{B_2}{s+1}+\frac{B_3}{s+2}
  \qquad (B_1,B_2,B_3\in\R^2)
\]
とおく．両辺に $s(s+1)(s+2)$ を掛けると
\[
  \binom{-4}{2s-2}
    =
    (s+1)(s+2)B_1
    +s(s+2)B_2
    +s(s+1)B_3.
\]
よって，$s=0,-1,-2$ を順に代入すれば
\begin{align*}
  B_1
    &=
    \left[
      \frac{1}{(s+1)(s+2)}
      \binom{-4}{2s-2}
    \right]_{s=0}
    =
    \binom{-2}{-1},\\
  B_2
    &=
    \left[
      \frac{1}{s(s+2)}
      \binom{-4}{2s-2}
    \right]_{s=-1}
    =
    \binom{4}{4},\\
  B_3
    &=
    \left[
      \frac{1}{s(s+1)}
      \binom{-4}{2s-2}
    \right]_{s=-2}
    =
    \binom{-2}{-3}.
\end{align*}
よって
\[
  \mathcal{L}^{-1}\!\left[(sE-A)^{-1}bU(s)\right]
    =
    \binom{-2}{-1}
    +
    e^{-t}\binom{4}{4}
    +
    e^{-2t}\binom{-2}{-3}.
\]
これは時間領域で直接積分した入力項と一致している．したがって，
式~\eqref{eq:solution-formula} で求めた解と同じ結果が得られる．
\end{proof}

\clearpage

\subsection{演習}

\begin{exe}
次の状態方程式の解を求め，また $t\to\infty$ のときの解の挙動を調べよ．
\[
  \begin{cases}
    \displaystyle
    \dot{x}
      =
      \begin{pmatrix}
        -4 & 1\\
        -6 & 1
      \end{pmatrix}x
      +
      \binom{1}{1}u,\\[3mm]
    u(t)=1.
  \end{cases}
\]
\end{exe}

\begin{proof}[解答]
\[
  A=
  \begin{pmatrix}
    -4 & 1\\
    -6 & 1
  \end{pmatrix},
  \qquad
  b=\binom{1}{1}
\]
とおく．まず $e^{tA}$ を求める．
\[
  \operatorname{tr}A=-4+1=-3,
  \qquad
  \det A=(-4)\cdot1-1\cdot(-6)=2
\]
であるから，特性方程式は
\[
  \lambda^2+3\lambda+2=0
  \qquad\Longleftrightarrow\qquad
  (\lambda+1)(\lambda+2)=0.
\]
よって固有値は $\lambda_1=-1,\lambda_2=-2$ である．相異なる実固有値の場合の公式より，
\[
  e^{tA}=e^{-t}P_1+e^{-2t}P_2
\]
である．ここで
\begin{align*}
  P_1
    &=
    \frac{1}{-1-(-2)}\{A-(-2)E\}
     =
    A+2E\\
    &=
    \begin{pmatrix}
      -4 & 1\\
      -6 & 1
    \end{pmatrix}
    +
    \begin{pmatrix}
      2 & 0\\
      0 & 2
    \end{pmatrix}
     =
    \begin{pmatrix}
      -2 & 1\\
      -6 & 3
    \end{pmatrix},\\
  P_2
    &=
    \frac{1}{-2-(-1)}\{A-(-1)E\}
     =
    -(A+E)\\
    &=
    -
    \left\{
    \begin{pmatrix}
      -4 & 1\\
      -6 & 1
    \end{pmatrix}
    +
    \begin{pmatrix}
      1 & 0\\
      0 & 1
    \end{pmatrix}
    \right\}
     =
    \begin{pmatrix}
      3 & -1\\
      6 & -2
    \end{pmatrix}.
\end{align*}
したがって
\[
  e^{tA}
    =
    e^{-t}
    \begin{pmatrix}
      -2 & 1\\
      -6 & 3
    \end{pmatrix}
    +
    e^{-2t}
    \begin{pmatrix}
      3 & -1\\
      6 & -2
    \end{pmatrix}.
\]

次に入力による項を計算する．ここでは $u(t)=1$ なので $bu(t)=b$ である．
$r=t-\tau$ とおくと，
\begin{align*}
  e^{rA}b
    &=
    e^{-r}
    \begin{pmatrix}
      -2 & 1\\
      -6 & 3
    \end{pmatrix}
    \binom{1}{1}
    +
    e^{-2r}
    \begin{pmatrix}
      3 & -1\\
      6 & -2
    \end{pmatrix}
    \binom{1}{1}\\
    &=
    e^{-r}\binom{-1}{-3}
    +
    e^{-2r}\binom{2}{4}.
\end{align*}
よって
\begin{align*}
  \int_0^t e^{(t-\tau)A}bu(\tau)\,d\tau
    &=
    \int_0^t
      \left\{
        e^{-r}\binom{-1}{-3}
        +
        e^{-2r}\binom{2}{4}
      \right\}dr\\
    &=
    \binom{-1}{-3}(1-e^{-t})
      +
      \binom{2}{4}\frac{1-e^{-2t}}{2}\\
    &=
    \binom{0}{-1}
      +
      e^{-t}\binom{1}{3}
      +
      e^{-2t}\binom{-1}{-2}.
\end{align*}
したがって，解の公式~\eqref{eq:solution-formula} より
\[
  x(t)
    =
    \left\{
      e^{-t}
      \begin{pmatrix}
        -2 & 1\\
        -6 & 3
      \end{pmatrix}
      +
      e^{-2t}
      \begin{pmatrix}
        3 & -1\\
        6 & -2
      \end{pmatrix}
    \right\}x(0)
    +
    e^{-t}\binom{1}{3}
    +
    e^{-2t}\binom{-1}{-2}
    +
    \binom{0}{-1}.
\]
右辺の $e^{-t}$，$e^{-2t}$ を含む項は $t\to\infty$ で消える．
したがって，初期値に依存しない定常値として
\[
  \lim_{t\to\infty}x(t)=\binom{0}{-1}
\]
である．
\end{proof}

\subsection{演習のラプラス変換による解法}

\begin{proof}[別解]
先ほどの演習
\[
  \dot{x}
    =
    \begin{pmatrix}
      -4 & 1\\
      -6 & 1
    \end{pmatrix}x
    +
    \binom{1}{1}u,
  \qquad
  u(t)=1
\]
をラプラス変換で解く．
\[
  A=
  \begin{pmatrix}
    -4 & 1\\
    -6 & 1
  \end{pmatrix},
  \qquad
  b=\binom{1}{1}
\]
とおくと，
\[
  sE-A
    =
    \begin{pmatrix}
      s+4 & -1\\
      6 & s-1
    \end{pmatrix}
\]
であり，
\[
  \det(sE-A)
    =(s+4)(s-1)+6
    =s^2+3s+2
    =(s+1)(s+2).
\]
したがって
\[
  (sE-A)^{-1}
    =
    \frac{1}{(s+1)(s+2)}
    \begin{pmatrix}
      s-1 & 1\\
      -6 & s+4
    \end{pmatrix}.
\]
これを
\[
  \frac{A_1}{s+1}+\frac{A_2}{s+2}
\]
と部分分数分解する．両辺に $(s+1)(s+2)$ を掛けると
\[
  \begin{pmatrix}
    s-1 & 1\\
    -6 & s+4
  \end{pmatrix}
    =
    (s+2)A_1+(s+1)A_2
\]
となる．よって，$s=-1,-2$ を代入して
\begin{align*}
  A_1
    &=
    \left[
      \frac{1}{s+2}
      \begin{pmatrix}
        s-1 & 1\\
        -6 & s+4
      \end{pmatrix}
    \right]_{s=-1}
     =
    \begin{pmatrix}
      -2 & 1\\
      -6 & 3
    \end{pmatrix},\\
  A_2
    &=
    \left[
      \frac{1}{s+1}
      \begin{pmatrix}
        s-1 & 1\\
        -6 & s+4
      \end{pmatrix}
    \right]_{s=-2}
     =
    \begin{pmatrix}
      3 & -1\\
      6 & -2
    \end{pmatrix}.
\end{align*}
したがって
\[
  \mathcal{L}^{-1}\!\left[(sE-A)^{-1}\right]
    =
    e^{-t}
    \begin{pmatrix}
      -2 & 1\\
      -6 & 3
    \end{pmatrix}
    +
    e^{-2t}
    \begin{pmatrix}
      3 & -1\\
      6 & -2
    \end{pmatrix}.
\]

また，$u(t)=1$ なので $U(s)=1/s$ である．入力による項は
\[
  (sE-A)^{-1}bU(s)
    =
    \frac{1}{(s+1)(s+2)}
    \begin{pmatrix}
      s-1 & 1\\
      -6 & s+4
    \end{pmatrix}
    \binom{1}{1}\frac{1}{s}
    =
    \frac{1}{s(s+1)(s+2)}
    \binom{s}{s-2}.
\]
これを
\[
  \frac{B_1}{s}+\frac{B_2}{s+1}+\frac{B_3}{s+2}
  \qquad (B_1,B_2,B_3\in\R^2)
\]
とおく．両辺に $s(s+1)(s+2)$ を掛けると
\[
  \binom{s}{s-2}
    =
    (s+1)(s+2)B_1
    +s(s+2)B_2
    +s(s+1)B_3.
\]
よって，$s=0,-1,-2$ を順に代入して
\begin{align*}
  B_1
    &=
    \left[
      \frac{1}{(s+1)(s+2)}
      \binom{s}{s-2}
    \right]_{s=0}
     =
    \binom{0}{-1},\\
  B_2
    &=
    \left[
      \frac{1}{s(s+2)}
      \binom{s}{s-2}
    \right]_{s=-1}
     =
    \binom{1}{3},\\
  B_3
    &=
    \left[
      \frac{1}{s(s+1)}
      \binom{s}{s-2}
    \right]_{s=-2}
     =
    \binom{-1}{-2}.
\end{align*}
したがって
\[
  \mathcal{L}^{-1}\!\left[(sE-A)^{-1}bU(s)\right]
    =
    \binom{0}{-1}
    +
    e^{-t}\binom{1}{3}
    +
    e^{-2t}\binom{-1}{-2}.
\]
以上より
\[
  x(t)
    =
    \left\{
      e^{-t}
      \begin{pmatrix}
        -2 & 1\\
        -6 & 3
      \end{pmatrix}
      +
      e^{-2t}
      \begin{pmatrix}
        3 & -1\\
        6 & -2
      \end{pmatrix}
    \right\}x(0)
    +
    e^{-t}\binom{1}{3}
    +
    e^{-2t}\binom{-1}{-2}
    +
    \binom{0}{-1}.
\]
これは時間領域で直接積分した解と一致する．また，指数関数を含む項は
$t\to\infty$ で消えるので
\[
  \lim_{t\to\infty}x(t)=\binom{0}{-1}
\]
である．
\end{proof}

\end{document}
